\documentclass[11pt,a4paper]{article}

% Essential packages
\usepackage[utf8]{inputenc}
\usepackage[margin=1in]{geometry}
\usepackage{amsmath,amssymb,amsthm}
\usepackage{graphicx}
\usepackage{booktabs}
\usepackage{hyperref}
\usepackage{caption}
\usepackage{subcaption}
\usepackage{float}

% Title and author
\title{[EXPERIMENT TITLE]\\
\large [Subtitle or Method Comparison]}
\author{[Author Names]}
\date{[Date]}

\begin{document}

\maketitle

\begin{abstract}
[One paragraph: hypothesis, method, key result, implication]
\end{abstract}

\section{Introduction}
[Motivation: why is this experiment important?]

[Problem statement: what gap are we addressing?]

[Hypothesis: what do we expect and why?]

\section{Methodology}

\subsection{Model Architecture}
[Describe the model/method being evaluated]

\begin{equation}
[Key equation 1: model formulation]
\label{eq:model}
\end{equation}

\subsection{Training Procedure}
[Describe training algorithm, optimizer, etc.]

\begin{equation}
[Key equation 2: loss or objective]
\label{eq:objective}
\end{equation}

\subsection{Experimental Setup}
[Dataset, splits, evaluation metrics]

\begin{table}[h]
\centering
\caption{Hyperparameters and Training Configuration}
\label{tab:hyperparameters}
\begin{tabular}{ll}
\toprule
Parameter & Value \\
\midrule
Base Model & [model name] \\
Learning Rate & [value] \\
Batch Size & [value] \\
Training Duration & [hours/iterations] \\
Hardware & [GPU specs] \\
\bottomrule
\end{tabular}
\end{table}

\section{Results}

\subsection{Training Dynamics}
[Describe Figure 1: training curves, convergence behavior]

\begin{figure}[H]
\centering
\includegraphics[width=\textwidth]{[experiment]_comparison.png}
\caption{Training dynamics comparison between baseline and treatment.}
\label{fig:training}
\end{figure}

\subsection{Quantitative Comparison}
[Present detailed results in table]

\begin{table}[H]
\centering
\caption{Accuracy comparison across test sets}
\label{tab:results}
\begin{tabular}{lrrr}
\toprule
Test Set & Baseline & Treatment & $\Delta$ \\
\midrule
[test set 1] & [value] & [value] & [delta] \\
[test set 2] & [value] & [value] & [delta] \\
\midrule
Average & [avg] & [avg] & [avg delta] \\
\bottomrule
\end{tabular}
\end{table}

\subsection{Key Findings}
\begin{itemize}
\item [Finding 1: main metric improvement]
\item [Finding 2: secondary observations]
\item [Finding 3: failure cases or limitations]
\end{itemize}

\begin{figure}[H]
\centering
\includegraphics[width=\textwidth]{[experiment]_key_metrics.png}
\caption{Detailed analysis: improvements, distributions, and pairwise comparison.}
\label{fig:detailed}
\end{figure}

\section{Analysis}
[Interpret the results: why did the treatment work or not work?]

[Address the hypothesis: supported or rejected?]

[Discuss implications: what do we learn about the method/problem?]

\section{Conclusion}
[Summarize: hypothesis status, key findings, impact]

[Future work: what experiments should follow?]

\section*{Acknowledgments}
[Optional: funding, compute resources, collaborators]

\begin{thebibliography}{9}
[Optional: key references]
\end{thebibliography}

\appendix
\section{Additional Details}
[Optional: full hyperparameter list, additional test sets, extended analysis]

\end{document}
